% Chapter 16 -- A small CPU that runs real programs
\chapter{A small CPU that runs real programs}\label{ch:cpu8}

Ch14 gave us an ALU that computes; ch15 gave it registers to remember
with and buses to share the wires. What is missing is the umpire: a
decoder that reads each instruction byte and drives the datapath's
control lines so the right thing happens at the right edge. Add that
and the pile of parts becomes a computer --- a real one, small but
complete, that assembles and runs real programs: it counts, it
multiplies, and it prints the Fibonacci sequence byte by byte on a
serial line. This chapter builds the C1 machine's last block, then
steps back to ask what we actually made --- and what we skipped on
the way to RISC-V.

% from-zero:begin
\begin{rulebox}[Before you start]
Chapters~14--15. \emph{Recap:} an instruction set architecture (ISA) is the contract between hardware and software: the list of instructions the CPU understands and what each one does.
\end{rulebox}
\begin{wins}[After this chapter you can]
\begin{itemize}
\item read a frozen ISA specification;
\item implement fetch, decode and execute for it in Verilog;
\item write a small assembler;
\item run a real program on your CPU. \textbf{Book~7} starts from this CPU and maps it to RISC-V.
\end{itemize}
\end{wins}
\begin{principle}[How to think about this]
A CPU is a loop: fetch an instruction, work out what it means, do it, repeat. The instruction set is the contract between the programmer and the machine, so freeze it before writing the hardware. Once your own CPU runs a program, you understand what every computer is doing underneath.
\end{principle}
% from-zero:end

\section{The ISA, frozen}\label{sec:isa16}

\index{instruction set architecture}\index{ISA}
Every CPU design starts with a contract: what instructions exist,
what each one costs, and how they are encoded. Ours is ten opcodes
over two fields --- the byte ch15 seeded, plus one honest upgrade.

\begin{figure}[htbp]
\centering
\begin{tikzpicture}[font=\scriptsize]
\node[anchor=north west, inner sep=0] at (0,0) {%
\begin{tabular}{lllll}
\toprule
opcode & mnemonic & operand & bytes & effect\\
\midrule
\texttt{0x0} & NOP   & ---     & 1 & nothing\\
\texttt{0x1} & LDI   & ---     & 2 & \texttt{A $\leftarrow$ imm8}\\
\texttt{0x2} & MOV   & \texttt{rd}    & 1 & \texttt{r[d] $\leftarrow$ A}\\
\texttt{0x3} & ADD   & \texttt{rs}    & 1 & \texttt{A $\leftarrow$ A + r[s]}, flags\\
\texttt{0x4} & SUB   & \texttt{rs}    & 1 & \texttt{A $\leftarrow$ A $-$ r[s]}, flags\\
\texttt{0x5} & AND   & \texttt{rs}    & 1 & \texttt{A $\leftarrow$ A \& r[s]}, flags\\
\texttt{0x6} & JMP   & ---     & 2 & \texttt{pc $\leftarrow$ imm8}\\
\texttt{0x7} & JZ    & ---     & 2 & \texttt{pc $\leftarrow$ imm8} if Z\\
\texttt{0x8} & OUT   & ---     & 1 & \texttt{out $\leftarrow$ A}\\
\texttt{0x9} & HALT  & ---     & 1 & pc holds forever\\
\texttt{0xA} & MOVA  & \texttt{rs}    & 1 & \texttt{A $\leftarrow$ r[s]}\\
\bottomrule
\end{tabular}};
\node[anchor=north west, text width=7.2cm, font=\scriptsize, text=slate] at (0,-6.2)
{Byte 0 is \texttt{\{op[7:4], operand[3:0]\}}; the operand nibble names
registers \texttt{r0}--\texttt{r7} (r0 doubles as the accumulator
\texttt{A}). Two-byte instructions carry an immediate byte at
\texttt{pc+1} --- including jump targets, the one place ch15's seeded
encoding needed a wider field: a 4-bit target cannot reach past
\texttt{0x0F}.};
\end{tikzpicture}

\caption{The C1 instruction set. Eleven words of vocabulary is all a
computer needs to be one. Regenerate: \ttfamily
figures/ch16/fig\_isa.tex.}
\label{fig:isa16}
\end{figure}

The upgrade deserves a sentence of its own. ch15 proposed jump
targets in the operand nibble --- four bits, sixteen addresses. The
multiply program alone is twenty-five bytes: a 4-bit jump cannot
reach most of it. So \texttt{JMP} and \texttt{JZ} became two-byte
instructions whose second byte is the whole target. The rule that
survives: \emph{any field that addresses the program wants as many
bits as the program space}, and our program space is now 256 bytes.
Everything else --- register fields, the ALU contract, the control
word --- is exactly as ch14 and ch15 left it.

\section{Fetch, decode, execute --- in one cycle}\label{sec:fde16}

\index{fetch-decode-execute}
Textbooks draw fetch, decode and execute as three phases. Ours all
happen inside one clock edge, and it is worth seeing why that is not
cheating:

\begin{enumerate}
\item \textbf{Fetch} is \texttt{rom[pc]} --- combinational read, so
the instruction is present all cycle.
\item \textbf{Decode} is a \texttt{case} statement on
\texttt{instr[7:4]} --- combinational, so the control word is present
all cycle.
\item \textbf{Execute} is the datapath settling --- combinational ---
then the edge: registers latch, flags latch, pc advances.
\end{enumerate}

One cycle per instruction, start to finish. The catch is the one we
flagged in ch15: the combinational path from \texttt{pc} through the
ROM, the decoder, the regfile, the ALU and the write mux into the
register file's data input is long. Real CPUs fix it by splitting the
cycle into stages --- which is exactly what a pipeline is, and why
this design will never win a frequency contest (Section
\ref{sec:missing16} is honest about that).

\begin{figure}[htbp]
\centering
\begin{tikzpicture}[font=\scriptsize, >=latex, scale=0.9, transform shape]
% datapath (left, shaded as the ch15 box)
\draw[thick, fill=graybg] (0,0.6) rectangle (7.2,6.2);
\node[font=\small\bfseries] at (3.6,5.9) {datapath (ch15)};
\draw[thick, fill=white] (0.5,4.4) rectangle (2.0,5.4);
\node[align=center] at (1.25,4.9) {PC /8};
\draw[thick, fill=white] (3.0,4.4) rectangle (5.0,5.4);
\node[align=center] at (4.0,4.9) {ROM 256x8};
\draw[->, thick] (2.0,4.9) -- (3.0,4.9) node[midway,above]{\texttt{pc}};
\draw[thick, fill=white] (0.5,1.6) rectangle (2.9,3.4);
\node[align=center] at (1.7,2.9) {regfile 8x8\\1W+2R};
\draw[thick, fill=accentsoft] (3.6,1.6) rectangle (6.0,3.0);
\node[align=center] at (4.8,2.6) {\texttt{alu8}};
\draw[->, thick] (2.9,2.5) -- (3.6,2.2) node[pos=0.5,below]{\texttt{A}};
\draw[->, thick] (2.9,2.0) -- (3.2,1.9) -- (4.2,1.9) -- (4.2,1.6);
\draw[thick, fill=goodbg] (5.0,3.6) rectangle (6.4,4.3);
\node[align=center] at (5.7,3.95) {flags};
\draw[->, thick] (4.8,3.0) -- (5.7,3.6);
\draw[thick, fill=goodbg] (6.2,1.0) rectangle (7.0,1.9);
\node[align=center, font=\tiny] at (6.6,1.45) {OUT};
\draw[->, thick] (6.6,1.9) -- (6.6,2.4) -- (6.0,2.4);
% decoder (right)
\draw[thick, fill=amberbg] (8.0,1.4) rectangle (10.6,5.6);
\node[align=center, font=\small\bfseries] at (9.3,5.2) {decoder};
\node[align=left] at (9.3,3.4) {\texttt{op[7:4] + flags}\\
$\to$ \texttt{alu\_op, wb\_sel,}\\
\texttt{dst\_sel, we,}\\
\texttt{flags\_we, out\_we,}\\
\texttt{pc\_mode}};
\draw[->, thick] (5.0,5.2) -- (5.6,5.2) -- (5.6,5.9) -- (8.8,5.9) -- (8.8,5.6);
\node[above] at (7.0,5.9) {\texttt{instr}};
\foreach \x/\y in {1.25/4.4, 1.7/1.6, 5.7/3.6, 6.6/1.0}{
  \draw[<-, dashed, slate] (\x,\y) -- (\x,\y-0.6);}
\foreach \x/\y in {1.25/4.4, 1.7/3.4, 4.8/3.0, 5.7/4.3, 6.6/1.9}{
  \draw[-, dashed, slate] (\x,\y) -- (8.0,\y);}
\draw[->, thick] (9.3,1.4) -- (9.3,0.4) -- (1.25,0.4) -- (1.25,0.6);
\node[align=center, text=slate, text width=10.4cm] at (5.1,-0.6)
{The centrepiece: every dashed line is a control-word tap. The decoder
is combinational --- fetch, decode and execute all happen inside one
cycle, which is why our ``master FSM'' has only two states.};
\end{tikzpicture}

\caption{The complete C1 machine --- the book's centrepiece. The ch15
datapath on the left gains a decoder on the right; the dashed taps
are the same 11-bit control word, now driven by hardware instead of
by a testbench. Regenerate: \ttfamily figures/ch16/fig\_cpu.tex.}
\label{fig:cpu16}
\end{figure}

\begin{figure}[htbp]
\centering
% traffic FSM state diagram (Moore)
\begin{tikzpicture}[font=\footnotesize, state/.style={draw, circle, minimum size=1.9cm, align=center, inner sep=1pt}]
\node[state] (g1) at (0,0) {NS\_G\\ \scriptsize NS=grn, EW=red};
\node[state] (y1) at (3.4,0) {NS\_Y\\ \scriptsize NS=ylw, EW=red};
\node[state] (g2) at (3.4,-3.0) {EW\_G\\ \scriptsize EW=grn, NS=red};
\node[state] (y2) at (0,-3.0) {EW\_Y\\ \scriptsize EW=ylw, NS=red};
\draw[->, thick] (g1) -- node[above] {timer} (y1);
\draw[->, thick] (y1) -- node[right] {timer} (g2);
\draw[->, thick] (g2) -- node[below] {timer} (y2);
\draw[->, thick] (y2) -- node[left] {timer} (g1);
\draw[->, dashed, ink] (-1.4,1.4) node[left] {rst} -- (g1.north west);
\end{tikzpicture}

\caption{The master FSM --- two states, deliberately small. RUN does
one whole instruction per edge because fetch/decode/execute are
combinational; HALT parks the pc. Regenerate: \ttfamily
figures/ch16/fig\_fsm.tex.}
\label{fig:fsm16}
\end{figure}

\section{The decoder: a case statement with a job title}\label{sec:decoder16}

\index{control unit}\index{decoder}
Here is the entire control unit, from \texttt{code/ch16\_cpu8.v}:

\begin{lstlisting}[language=Verilog, caption={The decoder: instruction
opcode in, 11-bit control word out --- combinational, so the whole
instruction completes on the next edge. From \ttfamily
code/ch16\_cpu8.v.}, label={lst:decoder16}]
always @(*) begin
  rf_we = 0; flags_we = 0; out_we = 0;
  wb_sel = 2'b00; dst_sel = 0; pc_mode = 2'b00; alu_op = 3'b000;
  case (op)
    4'h1: begin wb_sel=2'b01; rf_we=1; pc_mode=2'b01; end // LDI
    4'h2: begin wb_sel=2'b10; dst_sel=1; rf_we=1; end      // MOV rd
    4'h3: begin alu_op=3'b000; rf_we=1; flags_we=1; end    // ADD
    4'h4: begin alu_op=3'b001; rf_we=1; flags_we=1; end    // SUB
    4'h5: begin alu_op=3'b010; rf_we=1; flags_we=1; end    // AND
    4'h6: pc_mode = 2'b10;                                // JMP
    4'h7: pc_mode = flags[3] ? 2'b10 : 2'b01;             // JZ
    4'h8: out_we = 1;                                     // OUT
    4'h9: pc_mode = 2'b11;                                // HALT
    4'hA: begin wb_sel=2'b11; rf_we=1; end                // MOVA
    default: ;                                            // NOP
  endcase
end
\end{lstlisting}

Three things to read, slowly. \textbf{Every row is a control word}
--- match it against the ch15 table line by line; the hardware writes
the same rows the testbench played by hand. \textbf{Defaults first}:
every enable starts at 0, so an unrecognised opcode is a NOP, not a
crash --- ``silence is the safe state'' again. \textbf{\texttt{JZ}
reads the flags register}, not the ALU's live outputs: the Z flag was
latched by the previous arithmetic step, which is the only honest
answer ``did the last result equal zero.''

\begin{example}[Hand-trace before the machine]
Step through \texttt{A4 47 24} from the Fibonacci loop ---
\texttt{MOVA r4; SUB r7; MOV r4}:

\begin{center}
\begin{tabular}{llll}
\toprule
byte & decoded as & control word (key fields) & edge writes\\
\midrule
\texttt{a4} & \texttt{MOVA r4} & \texttt{wb=PB, dst=A, we=1, pc+1} & \texttt{r0$\leftarrow$r4}\\
\texttt{47} & \texttt{SUB r7} & \texttt{alu=001, we=1, fw=1, pc+1} & \texttt{r0$\leftarrow$A--r7}, Z NCV\\
\texttt{24} & \texttt{MOV r4} & \texttt{wb=PA, dst=r4, we=1, pc+1} & \texttt{r4$\leftarrow$A}\\
\bottomrule
\end{tabular}
\end{center}

Three cycles: load counter into A, decrement it (the SUB latches Z
--- the flag JZ consults next), store it back. The
accumulator-bottleneck rhythm is visible: every operation must pass
through A, which is exactly what a real instruction set costs when
the encoding budget is one byte.
\end{example}

\begin{checkpoint}
Why does the flags register exist at all? Try removing it: the JZ
step reads \texttt{alu\_z} live --- but by then the ALU inputs have
moved on. Flags must be \emph{latched by} the arithmetic step, not
recomputed by the branch step. That is the whole reason.
\end{checkpoint}

\section{The assembler}\label{sec:asm16}

\index{assembler}
Writing hex bytes by hand is how you learn encoding; writing a
60-line Python script is how you stop. \texttt{code/ch16\_asm8.py}
does two passes: first it walks the file counting bytes to give each
\texttt{label:} an address, then it emits \texttt{\{op,operand\}}
nibbles and immediate bytes.

\begin{lstlisting}[language=Python, caption={The assembler's core ---
from \ttfamily code/ch16\_asm8.py. Two passes: labels first, bytes
second.}, label={lst:asm16}]
OPS = {'NOP':0x0,'LDI':0x1,'MOV':0x2,'ADD':0x3,'SUB':0x4,'AND':0x5,
       'JMP':0x6,'JZ':0x7,'OUT':0x8,'HALT':0x9,'MOVA':0xA}
TWO = {'LDI','JMP','JZ'}                # ops carrying an imm byte
# pass 1: pc counts 2 bytes for TWO-ops, 1 otherwise -> label addr
# pass 2: out = (op<<4)|operand, then imm byte if needed
\end{lstlisting}

\begin{example}[Program 1: count on the LEDs]
\texttt{ch16\_count.asm} assembles to nine bytes:

\begin{lstlisting}[caption={\ttfamily ch16\_count.asm and its
assembly --- \ttfamily ch16\_count.hex.}]
LDI A,1 ; MOV r7        -> 10 01 27
LDI A,0 ; A = 0         -> 10 00
loop: ADD r7 ; A += 1   -> 37
      OUT               -> 80
      JMP loop          -> 60 06
\end{lstlisting}

Each lap: A grows by 1, OUT shows it, jump back to byte 6. On the
board this is the visible demo --- OUT drives the LED bank and the
UART stream.
\end{example}

\begin{example}[Program 2: multiply by repeated addition]
\texttt{ch16\_mult.asm} computes $7\times6$: a register counts down
from 6 while a second register accumulates 7 each lap. The loop is
the same decrement-jump-if-zero idiom as Fibonacci --- the
\texttt{SUB} both computes and sets the flag the \texttt{JZ} reads.

The bench runs it end to end. Real output, trimmed:

\begin{lstlisting}[caption={\ttfamily ch16\_cpu8\_mult\_tb.v running
\texttt{mult.hex} --- the machine prints 42 and parks.}]
PASS: 7*6 printed 42
mult run halted at pc=24
\end{lstlisting}
\end{example}

\begin{example}[Program 3: Fibonacci, byte by byte]
\texttt{ch16\_fib.asm} (the listing lives in \texttt{code/}) prints
each term mod 256. The testbench decodes every OUT write and checks
the whole sequence --- the real run:

\begin{lstlisting}[caption={\ttfamily ch16\_cpu8\_tb.v running
\texttt{fib.hex} --- twelve correct OUT bytes, then HALT.}]
OUT   1 ok (pc=14)      OUT  13 ok (pc=14)
OUT   1 ok (pc=14)      OUT  21 ok (pc=14)
OUT   2 ok (pc=14)      OUT  34 ok (pc=14)
OUT   3 ok (pc=14)      OUT  55 ok (pc=14)
OUT   5 ok (pc=14)      OUT  89 ok (pc=14)
OUT   8 ok (pc=14)      OUT 144 ok (pc=14)
PASS: cpu8 printed Fibonacci 1..144 then halted (12 OUT bytes)
ALL CHECKS PASS
\end{lstlisting}

Twelve instructions executed per term, roughly sixty cycles of real
hardware work --- and the bench verified every byte against the
expected sequence, not just the last one.
\end{example}

\section{Out the door: LEDs and UART}\label{sec:out16}

\index{UART}
The OUT register is where the program meets the world. The board
top file (\texttt{code/}) feeds the six onboard
directly, and every OUT write also launches a UART frame through the
BL702 bridge chip --- so the same byte appears on the LED bank and in
a serial monitor. \texttt{ch16\_uart\_tx8.v} is the smallest honest
transmitter: a four-state FSM shifting out a start bit, eight data
bits LSB-first, and a stop bit, one bit per baud interval.

\begin{figure}[htbp]
\centering
\begin{tikzpicture}[font=\scriptsize]
\begin{axis}[
  width=0.9\linewidth, height=4.6cm,
  xlabel={time (ns, simulation)}, ylabel={\texttt{tx}},
  xmin=0, xmax=800000, ymin=-0.4, ymax=1.5,
  ytick={0,1}, yticklabels={0,1},
  axis lines=left, clip=true]
\addplot[const plot, thick, accent] table[x=t, y=tx] {figures/ch16/uart.dat};
\draw[<-, slate] (axis cs:65000,0) -- (axis cs:65000,-0.3)
  node[below]{start};
\draw[<-, slate] (axis cs:225000,1) -- (axis cs:225000,1.42)
  node[above, font=\tiny]{\texttt{1 0 0 1 1 1 0 0} = 0x39, LSB first};
\draw[<-, slate] (axis cs:715000,1) -- (axis cs:715000,1.42)
  node[above]{stop};
\end{axis}
\end{tikzpicture}

\caption{A real UART frame on the wire --- the testbench's own VCD
output for byte \texttt{0x39}: start low, eight data bits LSB-first,
stop high. The companion bench decodes the same line the way a real
receiver would (wait for the start edge, sample mid-bit) and passes.
Regenerate data: \ttfamily figures/ch16/uart.dat.}
\label{fig:uart16}
\end{figure}

\begin{figure}[htbp]
\centering
\begin{tikzpicture}[font=\scriptsize, >=latex]
% board outline
\draw[thick, fill=graybg, rounded corners] (0,0) rectangle (8.8,4.6);
\node[font=\small\bfseries] at (4.4,4.3) {Tang Nano 9K --- cpu8 build};
% fpga
\draw[thick, fill=accentsoft] (3.4,1.7) rectangle (5.4,3.1);
\node[align=center] at (4.4,2.4) {GW1NR-LV9\\\texttt{cpu8}};
% leds
\foreach \i in {0,...,5}{
  \draw[thick, fill=goodbg] (6.4+\i*0.38,3.5) rectangle (6.7+\i*0.38,3.8);}
\node[below] at (7.3,3.4) {LEDs 0--5 = \texttt{out\_reg} (active low)};
\draw[->, slate] (5.4,2.8) -- (6.6,3.5);
% button
\draw[thick, fill=amberbg] (0.6,3.4) rectangle (1.6,3.9);
\node at (1.1,3.65) {BTN};
\node[below] at (1.1,3.3) {reset (pin 4)};
\draw[->, slate] (1.6,3.4) -- (3.4,2.6);
% uart
\draw[thick, fill=white] (0.6,1.0) rectangle (2.6,2.0);
\node[align=center] at (1.6,1.5) {BL702\\USB-UART};
\draw[->, slate] (2.6,1.5) -- (3.4,2.0) node[midway, below]{\texttt{tx} pin 17};
% osc
\draw[thick, fill=white] (6.4,0.8) rectangle (7.4,1.6);
\node[align=center] at (6.9,1.2) {27 MHz};
\draw[->, slate] (6.9,1.6) -- (4.9,2.1) node[midway, below]{\texttt{clk} pin 52};
% usb
\draw[thick, fill=white] (0.6,0.2) rectangle (2.6,0.7);
\node at (1.6,0.45) {USB-C};
\node[align=left, text=slate, text width=10cm] at (5.0,-0.7)
{Wiring it runs: \texttt{i\_clk27} 52, \texttt{i\_btn} 4,
\texttt{o\_led[0..5]} 10--16, \texttt{o\_uart\_tx} 17 (the BL702
bridge: open a serial monitor at 115200 and the OUT bytes arrive).};
\end{tikzpicture}

\caption{The C1 machine on the Tang Nano 9K: clock, reset button, LED
bank, and the UART bridge that carries OUT bytes to a serial monitor.
Regenerate: \ttfamily figures/ch16/fig\_board16.tex.}
\label{fig:board16}
\end{figure}

\begin{example}[What the whole machine costs]
The real \texttt{synth\_gowin} stat for \texttt{cpu8}, trimmed:

\begin{lstlisting}[caption={A working CPU in 302 cells --- less than
4\% of the smallest FPGA this book uses.}]
=== cpu8 ===
      302 cells
       30   ALU        8   DFF         1   DFFR        17   DFFRE
        4   RAM16SDP4  15   LUT1       21   LUT2       106   LUT3+
       51   MUX2_LUT*   2   IBUF       34   OBUF
\end{lstlisting}

Reading it against ch14's ALU stat: 30 ALU cells (the ALU plus the
pc adder), 26 flip-flops (pc8 + flags4 + out8 + divider bits), and
the register file mapped onto the same 4 distributed-RAM cells. An
entire stored-program computer for 302 cells --- this is why ``a CPU
is a finite state machine with very good bookkeeping'' is only half
a joke.
\end{example}

\textbf{Honest marking.} The Fibonacci run, the multiply run, the
count assembly, the UART frame decode and the synth stat are all real
outputs captured on this machine. The board build (\texttt{top\_cpu9k}
+ \texttt{ch16\_cpu9k.cst}) is complete and lints clean, but no board
was attached in this session: pin numbers follow the Sipeed rev-3718
schematic, and the clock divider brings the CPU to $\sim$6~Hz so the
LED count is visible --- that last step is verified by construction,
not by hardware.

\section{What this CPU is missing}\label{sec:missing16}

\index{CPU!limitations}
Being honest about the gaps is the point of building small:

\begin{itemize}
\item \textbf{No stack.} \texttt{CALL}/\texttt{RET} need somewhere to
put the return address; our pc has one mode, not two.
\item \textbf{No interrupts.} Nothing external can redirect the pc;
the machine is deaf until it HALTs.
\item \textbf{No pipeline.} One cycle per instruction means the
critical path is the whole machine. Figure~\ref{fig:pipe16} shows the
upgrade shape.
\item \textbf{No memory-mapped data.} Programs can only name the eight
registers --- there is no load/store into RAM at all.
\item \textbf{Immediate-in-rom only.} Constants live next to opcodes;
there is no data space to read them from.
\end{itemize}

\begin{figure}[htbp]
\centering
\begin{tikzpicture}[font=\scriptsize, >=latex]
% single cycle timeline
\node[font=\small\bfseries, anchor=west] at (0,4.6) {ours: one stage};
\draw[thick] (0,4.0) rectangle (6.4,4.4);
\node at (3.2,4.2) {fetch + decode + execute (1 cycle)};
\node[anchor=west, text=slate] at (6.6,4.2) {next instr waits};
% pipelined timeline
\node[font=\small\bfseries, anchor=west] at (0,3.0) {pipelined: three stages overlap};
\foreach \i/\s in {0/{i$_n$: fetch},2.2/{i$_n$: dec},4.4/{i$_n$: exe}}{
  \draw[thick, fill=accentsoft] (\i,2.2) rectangle (\i+1.8,2.8);
  \node at (\i+0.9,2.5) {\s};}
\foreach \i/\s in {2.2/{i$_{n+1}$: fetch},4.4/{i$_{n+1}$: dec},6.6/{i$_{n+1}$: exe}}{
  \draw[thick, fill=goodbg] (\i,1.4) rectangle (\i+1.8,2.0);
  \node at (\i+0.9,1.7) {\s};}
\foreach \i/\s in {4.4/{i$_{n+2}$: fetch},6.6/{i$_{n+2}$: dec},8.8/{i$_{n+2}$: exe}}{
  \draw[thick, fill=amberbg] (\i,0.6) rectangle (\i+1.8,1.2);
  \node at (\i+0.9,0.9) {\s};}
\node[align=left, text=slate, text width=10cm] at (5.4,-0.6)
{The teaser, not the lesson: overlap and throughput triple, but now the
registers between stages are real hardware and hazards appear ---
\texttt{ADD} reading \texttt{r1} while the previous \texttt{LDI} has not
written it yet. Book 7 opens here.};
\end{tikzpicture}

\caption{The pipeline teaser: overlap fetch/decode/execute across
three instructions and throughput triples --- but the latches between
stages and the hazards between operands are what Book 7 spends real
pages on. Regenerate: \ttfamily figures/ch16/fig\_pipe.tex.}
\label{fig:pipe16}
\end{figure}

\subsection{Where RISC-V picks up the thread}\label{sec:riscv16}

\index{RISC-V}
Every instruction we invented has a grown-up cousin in RV32I --- the
mapping is worth a paragraph because Book 7 lands as ``the same
ideas, done properly'':

\begin{center}
\begin{tabular}{lll}
\toprule
C1 & RV32I & the idea that survives\\
\midrule
\texttt{LDI A,imm} & \texttt{addi/lui} & constants ride in the instruction\\
\texttt{MOV / MOVA} & \texttt{addi rd,rs,0} & register moves are a degenerate add\\
\texttt{ADD/SUB/AND} & R-type ops & register$\to$register through the ALU\\
\texttt{JMP} & \texttt{jal x0} & unconditional redirect\\
\texttt{JZ} & \texttt{beq} & conditional branch on a flag-like test\\
\texttt{OUT} & store to a mmio port & I/O is just a write to an address\\
\texttt{HALT} & \texttt{ebreak} & handing the machine back\\
\bottomrule
\end{tabular}
\end{center}

The differences are the lessons: RISC-V replaces our lone accumulator
with 32 symmetric registers (the bottleneck dissolves), puts load and
store into a real data space (the RAM we simulated earns its keep),
and splits instructions across fixed fields big enough to never need
our operand-nibble compromises. Same machine, bigger budget.

\section{Common mistakes}\label{sec:cm16}

\begin{tripwire}[Encoding bugs the assembler accepts silently]
An assembler is a lookup table: \texttt{MOV r9} could emit a byte
that decodes to something else (or nothing) without a complaint.
Ours rejects out-of-range registers --- check yours does. The deeper
version of this trap: an ISA revision (like our jump-target widening)
that the assembler ``helpfully'' hides --- then sim and hardware
disagree and both look right.
\end{tripwire}
\begin{tripwire}[Program-counter edge cases]
\texttt{JMP} to yourself loops forever by design; \texttt{HALT}
mid-fetch must park the pc \emph{before} the next byte is read as an
opcode. And the reset question: our pc wakes at 0 --- if your program
doesn't start at 0, the first byte it executes is whatever the ROM
left there (usually \texttt{0x00}, a NOP --- harmless only by luck).
\end{tripwire}
\begin{tripwire}[A testbench that only ever runs one program]
The first version of this bench checked the final OUT byte --- it
passed while printing the \emph{previous} loop's value (the pre-edge
read: we caught it because the sequence check counted terms, not just
the last one). Check every output the program emits, in order ---
a single final check can pass a broken machine.
\end{tripwire}
\begin{tripwire}[Flags sampled at the wrong edge]
JZ reading \texttt{alu\_z} live fails: the ALU's inputs are already
the next instruction's operands. Flags belong to the operation that
produced them --- latch on \texttt{flags\_we}, read the register.
\end{tripwire}
\begin{tripwire}[CDC in the board glue]
The CPU at 6~Hz and the UART at 27~MHz are two clock domains; a
one-cycle pulse in the slow domain is hundreds of cycles wide to the
fast one (fine here --- we catch it level-style), but a fast-domain
pulse into a slow domain would need a real synchroniser. Ask which
direction the pulse travels before wiring domains together.
\end{tripwire}
\begin{tripwire}[Assuming unwritten state is zero]
The sim's \texttt{initial} blocks are honest in iverilog; real
flip-flops power up to vendor-defined or random states. We reset pc,
flags, OUT and the halted flag explicitly; the register file contents
are whatever they power up as --- programs must write before they
read (the ch15 lesson, now on hardware).
\end{tripwire}

\begin{exercise}[easy]
Why does \texttt{JZ} need to read the flags \emph{register} rather
than the ALU's live flag outputs?
\end{exercise}
\begin{exercise}[medium]
Trace bytes \texttt{60 06} back through the decoder: what is in each
control field during a JMP, and why is \texttt{wb\_sel} a don't-care?
\end{exercise}
\begin{exercise}[medium]
Our \texttt{MOV} is one-directional; we added \texttt{MOVA} rather
than redefining the operand nibble. List one alternative encoding
that would have worked and its cost.
\end{exercise}
\begin{exercise}[hard]
The decoder's defaults make unknown opcodes NOPs. Argue both sides:
is silently-NOP-ing bad opcodes a feature or a bug? What does
\texttt{ebreak} on RISC-V suggest?
\end{exercise}

\section{Summary}\label{sec:sum16}

\begin{center}
\begin{tabularx}{\linewidth}{@{}lX@{}}
\toprule
\textbf{Thing} & \textbf{The rule of thumb}\\
\midrule
ISA & the contract first: opcodes, fields, byte lengths --- freeze it
before the datapath\\
fetch-decode-exec & can live in one cycle when the datapath is
combinational --- the price is the critical path\\
control unit & a decoder is a case statement whose outputs are
write-enables and mux selects\\
flags & latch at compute time, read at branch time --- never recompute\\
assembler & 60 lines buys you labeled programs; two passes: addresses,
then bytes\\
verified run & check the whole output stream in order, not the last
byte\\
RISC-V bridge & every C1 instruction is a crippled RV32I one --- same
ideas, real fields\\
\bottomrule
\end{tabularx}
\end{center}

\section*{Additional exercises}
\addcontentsline{toc}{section}{Additional exercises}

\begin{exercise}[easy]
Count the C1 ISA: how many one-byte and two-byte instructions, and
what does the mix tell you about the encoding's pressure points?
\end{exercise}
\begin{exercise}[easy]
Why does the operand nibble suffice for registers but not for jump
targets?
\end{exercise}
\begin{exercise}[easy]
What state does the master FSM hold that the datapath doesn't?
\end{exercise}
\begin{exercise}[easy]
Name the two clock domains in \texttt{top\_cpu9k} and the signal that
crosses between them.
\end{exercise}
\begin{exercise}[medium]
Write \texttt{SUBI} (subtract immediate) end-to-end: encoding slot,
decoder row, assembler line, and a test program. Which existing
control fields change?
\end{exercise}
\begin{exercise}[medium]
The multiply program uses \texttt{SUB} to both decrement and set Z.
If \texttt{JZ} read a ``decrement flag'' instead, what breaks when the
counter hits 0x00 vs 0x01?
\end{exercise}
\begin{exercise}[medium]
The count program wraps at 255 silently. Write the check (register
contents + one added instruction) that makes wrap visible on the
LEDs.
\end{exercise}
\begin{exercise}[medium]
Our UART drops OUT bytes when busy. For the 6~Hz machine that's fine;
compute the minimum CPU rate at which drops become possible.
\end{exercise}
\begin{exercise}[hard]
Design a 3-bit ``opcode extension'' field so the ISA can grow past
0xF opcodes without touching instruction length. Where does the
decode table get its second dimension?
\end{exercise}
\begin{exercise}[hard]
A classmate pipelines our CPU by registering the decoder's outputs.
Find the hazard it introduces (hint: \texttt{JZ} and the flags reg),
and the cheapest fix.
\end{exercise}
\begin{exercise}[hard]
Write the truth table for \texttt{pc\_mode} including what must be
true when \texttt{JZ} evaluates. What invariant does
\texttt{flags[3]} encode?
\end{exercise}
\begin{exercise}[hard]
The UART testbench decodes mid-bit like a real RX. Extend it to check
stop-bit timing, then argue why \texttt{CLKS\_PER\_BIT=234} at 27 MHz
is correct for 115200 baud (or find the error).
\end{exercise}
\begin{exercise}[hard]
Port the Fibonacci program to the RV32I mapping in
\S\ref{sec:riscv16} --- write real RV32I instructions for each C1
line and count the cycle difference.
\end{exercise}
\begin{exercise}[hard]
Find the planted bug: a decoder row sets \texttt{pc\_mode=2'b00} on
\texttt{JZ}-taken instead of \texttt{2'b10}. What program behaviour
does it produce, and which bench check catches it first?
\end{exercise}
\begin{exercise}[hard]
Our register file is r0=accumulator, not r0=zero like RISC-V.
Estimate the code-size cost of the accumulator design: how many extra
bytes does the fib program pay versus a symmetric-register version?
\end{exercise}
\begin{exercise}[hard]
Design the \texttt{CALL}/\texttt{RET} pair: which new datapath
element is mandatory (not just nice), and sketch the two extra
pc\_mode rows.
\end{exercise}
\begin{exercise}[hard]
The synth stat shows 34 OBUF --- far more than the 6 LEDs + 1 TX pin
the board uses. Explain the discrepancy and what a real \texttt{.cst}
build would do with them.
\end{exercise}
\begin{exercise}[hard]
Prove the C1 machine is Turing-complete modulo memory size (or find
the instruction it lacks to be one). What does a jump-table entry
cost in our encoding?
\end{exercise}
\begin{exercise}[hard]
Write the multi-cycle variant of the decoder (FETCH and EXECUTE as
separate states): which control signals may stay combinational and
which must be registered? Sketch the new FSM.
\end{exercise}

\section*{Hints and answers}
\addcontentsline{toc}{section}{Hints and answers}
\begin{description}
\item[JZ and the flags register (easy, inline)] the live ALU flags
reflect whatever the ALU's inputs are \emph{now} --- the next
instruction's operands. The register freezes the Z computed by the
SUB that the branch actually asks about.
\item[JMP fields (medium)] op=6, so \texttt{pc\_mode=jump}, all
enables 0; \texttt{wb\_sel}/\texttt{alu\_op} irrelevant because
\texttt{we=0} --- nothing writes, so the datapath result dies on the
mux, which is the whole point of default-safe enables.
\item[One-nibble alternative (medium)] make operand[3] a direction
bit on MOV (0=write reg, 1=read into A); costs the ability to name
r4--r7 in one direction, and adds an asymmetry the assembler must
hide.
\item[Unknown-opcode policy (hard)] treating bad opcodes as NOP
hides encoding bugs (bad assemblies run ``cleanly''); trapping
surfaces them. The
\texttt{ebreak} on RISC-V exists to make ``stop here'' explicit ---
ours is friendlier to noise, theirs to correctness.
\item[1-vs-2 byte mix (easy)] 8 one-byte vs 3 two-byte (LDI, JMP,
JZ): immediates are the only thing too wide for a nibble.
\item[Nibble vs byte (easy)] 3 bits name 8 registers; a jump target
must reach all 256 program bytes.
\item[FSM state (easy)] only ``halted'' --- pc, flags and registers
are datapath state.
\item[Clock domains (easy)] the \texttt{div[21]} CPU clock vs the
27~MHz UART clock; \texttt{out\_we} and \texttt{out\_reg} cross.
\item[SUBI (medium)] new opcode 0xB, alu\_op=SUB, wb=ALU,
i\_b$\leftarrow$imm8 (the B-source mux is the only real datapath
change).
\item[Decrement flag (medium)] nothing breaks --- Z is Z either way;
the point is any flag-based branch must agree on which operation
last wrote flags.
\item[Visible wrap (medium)] check C flag after ADD: if carry-out,
increment a second OUT channel or latch an LED. Needs flags exposed
to the datapath --- they are, via the register.
\item[UART drop rate (medium)] a frame takes $10/115200$ seconds,
about 87~$\mu$s; safe while the OUT rate stays under
$1/87\mu s \approx 11.5$~kHz --- our 6~Hz is far inside budget.
\item[Extension field (hard)] a second nibble can't hold it ---
escape opcode 0xF to mean ``look at the next byte''; decode becomes
two-level and the assembler needs a prefix table.
\item[Pipelined decode hazard (hard)] the JZ row's pc\_mode depends on
flags that may still be in flight from the previous instruction ---
registering the decoder delays pc\_mode one cycle; fix: stall the
fetch when op is JZ (one bubble) or forward flags.
\item[pc\_mode truth (hard)] 00=+1, 01=+2, 10=imm8 jump, 11=hold;
JZ needs ``flags latched before decode reads them'' --- the
invariant is that flags are one-instruction-old, always.
\item[Baud check (hard)] $27\times10^6/234=115{,}384$~Hz --- within
0.2\% of 115200, standard tolerance; a real UART also checks the stop
bit arrives high one bit-period after the last data bit.
\item[RV32I port (hard)] \texttt{li a0,0; li a1,1};\newline
\texttt{loop: add a2,a1,a0; mv a0,a1; mv a1,a2; bnez} --- roughly
half the instructions: symmetric registers kill the A-bottleneck's
load/store shuffles.
\item[Planted JZ bug (hard)] taken jumps fall through to
\texttt{pc+1}: the \texttt{done:} branch is never taken ---
Fibonacci keeps printing past 144 until OUT 233+ deviates; the
sequence check fires first.
\item[Accumulator tax (hard)] every ``reg$\to$reg'' move costs a
round trip through A (2 instr) plus an LDI whenever A is clobbered
--- fib pays ~8 bytes vs a symmetric version; the tax scales with
how much data the program juggles.
\item[CALL/RET (hard)] a return-address stack (or a link register):
\texttt{CALL} pushes pc+2 then jumps; \texttt{RET} pops into pc ---
needs a stack pointer register and a memory write port the datapath
doesn't have. Two pc\_mode rows: push-then-jump, pop.
\item[34 OBUF (hard)] the debug taps: a8+pc8+instr8+out8+flags4 minus
... count the ports --- on a real build they cost nothing if left
unconnected, but top-level synthesis can't know that; strip debug
taps for the board build.
\item[Turing-completeness (hard)] conditional jump + read/write
registers + unbounded program space = a register machine (equivalent
to TM in principle); a jump table costs one byte per entry ---
fine at our scale.
\item[Multi-cycle decoder (hard)] FETCH latches instr (registered);
EXECUTE keeps control word combinational --- the split halves the
critical path; FSM gains FETCH/EXEC states and a 2-cycle CPI cost.
\end{description}

\section*{Further reading}
\addcontentsline{toc}{section}{Further reading}

\begin{itemize}
\item Nisan \& Schocken, \emph{The Elements of Computing Systems},
projects 4--6 (machine language, assembler, the Hack CPU).\newline
\url{nand2tetris.org/project05} --- the closest sibling to C1 in the
wild; our control-word table is their Hack control bits.
Intermediate. Verified 200, 2026-10-11.
\item Harris \& Harris, \emph{Digital Design and Computer
Architecture}, ch.~7 (microarchitecture) --- single-cycle datapath
and the control-FSM method, formalised.\newline
\url{pages.hmc.edu/harris/ddca/} --- Intermediate. Verified 200,
2026-10-11.
\item Coursera / Hebrew University, \emph{From Nand to Tetris,
Part~I} --- free course building the Hack CPU end to end.\newline
\url{coursera.org/learn/build-a-computer} --- Beginner. Verified
200, 2026-10-11.
\item Lushay Labs, \emph{Tang Nano series --- Our First CPU}.\newline
\url{learn.lushaylabs.com/tang-nano-series/} --- a different small
CPU on this exact board, with the Gowin toolchain steps spelled out.
Beginner--intermediate. Verified 200, 2026-10-11.
\item \emph{HDLBits}, Circuits/Finite State Machines.\newline
\url{hdlbits.01xz.net} --- the FSM drills behind our two-state
master. Beginner. Verified 200, 2026-10-11.
\item Wikipedia, \emph{Central processing unit} and
\emph{Instruction cycle}.\newline
\url{en.wikipedia.org/wiki/Central_processing_unit}\newline
\url{en.wikipedia.org/wiki/Instruction_cycle} --- the standard
vocabulary for what we built (and the real taxonomy of what we
skipped). Reference. Verified 200, 2026-10-11.
\end{itemize}

\begin{rulebox}[{Watch instead}]
\begin{enumerate}
\item \emph{8-bit CPU control logic: Part 1} --- Ben Eater, 27:52,
\url{https://www.youtube.com/watch?v=dXdoim96v5A} --- the control
unit built wire by wire on a breadboard: microcode, instruction
fetch, the step counter. Matches \S\ref{sec:decoder16} and our
decoder-vs-microcode discussion. Opened and confirmed 2026-10-11.
\item \emph{Instructions \& Programs: Crash Course Computer
Science \#8} --- CrashCourse, 10:36,
\url{https://www.youtube.com/watch?v=zltgXvg6r3k} --- the stored-program
idea and fetch-decode-execute as a concept, fast and visual. Matches
\S\ref{sec:fde16}. Opened and confirmed 2026-10-11.
\end{enumerate}
\end{rulebox}

\chapternotes{The ISA and single-cycle structure follow the
nand2tetris Hack conventions adapted to our register-accumulator
encoding; control-unit-as-decoder is Harris \& Harris's ch.~7
pattern \cite{ddca}; the UART transmitter and board constraints are
ours (Tang Nano 9K rev-3718 pin map \cite{apicula}); synth numbers
are real \texttt{yosys} 0.69 runs. The pipeline teaser and RV32I
mapping set up Book~7.}
